\begin{proof}
\textbf{Step 1. Degenerate / edge cases.}  

\begin{enumerate}
\item \textbf{Case \(n=1\).}  
Since \(\displaystyle\sum_{i=1}^{1}\lambda_{i}x_{i}=x_{1}\) and
\(\displaystyle\sum_{i=1}^{1}\lambda_{i}f(x_{i})=f(x_{1})\), inequality \((J)\) reads
\(f(x_{1})\le f(x_{1})\), which holds with equality.

\item \textbf{Zero weights.}  
If \(\lambda_{k}=0\) for some index \(k\), the term \(\lambda_{k}f(x_{k})\) vanishes
and the left–hand side does not contain \(x_{k}\). Hence the inequality for the
full family \(\{x_{i},\lambda_{i}\}\) is equivalent to the inequality for the
reduced family obtained by deleting all indices with zero weight. Therefore it
suffices to prove \((J)\) under the assumption
\[
\lambda_{i}>0 \quad\text{for all }i.
\]

\item \textbf{A weight equal to \(1\).}  
If \(\lambda_{j}=1\) (and consequently all other \(\lambda_{i}=0\)), then
\(\sum_{i}\lambda_{i}x_{i}=x_{j}\) and \(\sum_{i}\lambda_{i}f(x_{i})=f(x_{j})\);
again \((J)\) becomes an equality.
\end{enumerate}
Thus, apart from these trivial situations, we may assume
\[
n\ge 2,\qquad \lambda_{i}>0\;(i=1,\dots ,n). \tag{1}
\]

\bigskip

\textbf{Step 2. Base case \(n=2\).}  

For two points \(x_{1},x_{2}\) and a weight \(\lambda\in[0,1]\) (so the second
weight is \(1-\lambda\)), the definition of convexity gives
\begin{equation*}
f\bigl(\lambda x_{1}+(1-\lambda)x_{2}\bigr)\le
\lambda f(x_{1})+(1-\lambda)f(x_{2}). \tag{2.1}
\end{equation*}
Equality in \((2.1)\) holds \emph{iff} \(\lambda\in\{0,1\}\) or \(f\) is affine on
the segment \([x_{1},x_{2}]\) (in particular when \(x_{1}=x_{2}\)). Hence \((J)\)
is true for \(n=2\).

\bigskip

\textbf{Step 3. Inductive hypothesis.}  

Assume that for some integer \(m\) with \(2\le m\le n-1\) the inequality
\((J)\) holds for any collection of \(m\) points and any non‑negative weights
summing to \(1\); i.e.
\begin{equation*}
f\!\left(\sum_{i=1}^{m}\mu_{i}x_{i}\right)\le
\sum_{i=1}^{m}\mu_{i}f(x_{i})\qquad
\bigl(\mu_{i}\ge0,\;\sum_{i=1}^{m}\mu_{i}=1\bigr). \tag{3.1}
\end{equation*}

\bigskip

\textbf{Step 4. Inductive step \(m=n-1\Longrightarrow m=n\).}  

Let the weights \(\lambda_{1},\dots ,\lambda_{n}\) satisfy (1).  
The case \(\lambda_{1}=1\) is covered by Step 1.3, so we may assume
\[
0<\lambda_{1}<1 .
\]

Define
\[
y:=\frac{\displaystyle\sum_{i=2}^{n}\lambda_{i}x_{i}}{1-\lambda_{1}} .
\]
Because
\[
\sum_{i=2}^{n}\frac{\lambda_{i}}{1-\lambda_{1}}=1,
\]
the numbers
\[
\mu_{i}:=\frac{\lambda_{i}}{1-\lambda_{1}},\qquad i=2,\dots ,n,
\]
form a legitimate set of weights for the points \(x_{2},\dots ,x_{n}\).

\medskip

\textbf{4.1 Apply the inductive hypothesis to \(\{x_{2},\dots ,x_{n}\}\).}  
Using \((3.1)\) with the points \(x_{2},\dots ,x_{n}\) and the weights \(\mu_{i}\) we obtain
\begin{equation*}
f(y)=f\!\left(\sum_{i=2}^{n}\mu_{i}x_{i}\right)
\le \sum_{i=2}^{n}\mu_{i}f(x_{i})
 =\frac{1}{1-\lambda_{1}}\sum_{i=2}^{n}\lambda_{i}f(x_{i}). \tag{4.1}
\end{equation*}
Multiplying \((4.1)\) by the positive factor \(1-\lambda_{1}\) yields
\begin{equation*}
(1-\lambda_{1})\,f(y)\le\sum_{i=2}^{n}\lambda_{i}f(x_{i}). \tag{4.2}
\end{equation*}

\medskip

\textbf{4.2 Convexity step for the pair \((x_{1},y)\).}  
Applying the two‑point convexity \((C)\) with weight \(\lambda_{1}\) gives
\begin{equation*}
f\bigl(\lambda_{1}x_{1}+(1-\lambda_{1})y\bigr)
\le \lambda_{1}f(x_{1})+(1-\lambda_{1})f(y). \tag{4.3}
\end{equation*}

\medskip

\textbf{4.3 Identify the aggregated point.}  
By the definition of \(y\),
\begin{equation*}
\lambda_{1}x_{1}+(1-\lambda_{1})y
=\lambda_{1}x_{1}+\sum_{i=2}^{n}\lambda_{i}x_{i}
=\sum_{i=1}^{n}\lambda_{i}x_{i}. \tag{4.4}
\end{equation*}
Thus the left‑hand side of \((4.3)\) is exactly
\(f\!\left(\sum_{i=1}^{n}\lambda_{i}x_{i}\right)\).

\medskip

\textbf{4.4 Assemble the inequalities.}  
Insert \((4.2)\) into the right‑hand side of \((4.3)\) and use \((4.4)\):
\begin{align*}
f\!\left(\sum_{i=1}^{n}\lambda_{i}x_{i}\right)
&\le \lambda_{1}f(x_{1})+(1-\lambda_{1})f(y) \\
&\le \lambda_{1}f(x_{1})+\sum_{i=2}^{n}\lambda_{i}f(x_{i})
   =\sum_{i=1}^{n}\lambda_{i}f(x_{i}).
\end{align*}
Hence \((J)\) holds for \(n\) points, completing the induction.

\bigskip

\textbf{Step 5. Equality condition.}  

Equality in \((J)\) occurs precisely when equality holds in both \((4.3)\) and
\((4.2)\).

\begin{itemize}
\item \emph{Equality in \((4.3)\).}  
By the definition of convexity, equality holds iff either \(\lambda_{1}\in\{0,1\}\)
(the trivial weight case) or \(f\) is affine on the segment joining \(x_{1}\) and
\(y\) (in particular when \(x_{1}=y\)).

\item \emph{Equality in \((4.2)\).}  
Equality in the inductive hypothesis \((4.1)\) holds iff either one of the
weights \(\lambda_{2},\dots ,\lambda_{n}\) equals \(1\) (the rest being \(0\)),
or \(f\) is affine on the convex hull of \(\{x_{2},\dots ,x_{n}\}\) (in
particular when all those points coincide).
\end{itemize}

Combining these observations and recalling that the cases where a weight
equals \(1\) have already been dealt with in Step 1, we obtain the following
complete description:

\begin{enumerate}
\item If at least one weight equals \(1\) (and the others are zero), equality
holds (trivial case).
\item Otherwise all weights are strictly positive.  
Equality holds \emph{iff} the points with positive weight are all identical,
i.e.
\[
x_{1}=x_{2}= \dots =x_{n},
\]
or more generally iff \(f\) is affine on the convex hull of the points
\(\{x_{i}\}\). For a strictly convex function the only non‑trivial equality
case is the former.
\end{enumerate}

\bigskip

\textbf{Step 6. Conclusion.}  
We have proved Jensen’s inequality \((J)\) for any finite collection of real
numbers and any convex function, treated all degenerate situations explicitly,
and gave a rigorous derivation of the equality cases.

\qed
\end{proof}